\documentclass[10pt,a4paper]{amsart}
\usepackage[margin=19mm]{geometry}
\usepackage{amsmath,amssymb,mathtools}
\usepackage[scr=rsfs,cal=euler]{mathalfa}
\usepackage{xfrac}
\usepackage[unicode,hidelinks,bookmarks=false]{hyperref}
\newcommand{\ie}{i.e.\ }
\newcommand{\cf}{cf.\ }
\newcommand{\resp}{respectively\ }
\newcommand{\Subsection}{\subsection}
\providecommand{\nobreakdash}{\nobreak\hbox{-}}
\setlength{\emergencystretch}{2em}
\multlinegap0pt
\usepackage{bm}
\usepackage{amssymb}
\usepackage{mathtools}
\usepackage[shortlabels]{enumitem}
\usepackage{csquotes}
\newtheorem{lemma}{Lemma}
\title{A contact version of Kirby's theorem}
\author{Marc Kegel, Eric Stenhede,, Vera V\'ertesi}
\date{Source: arXiv:2605.27126v1 (CC BY 4.0)}
\begin{document}
\maketitle

\noindent\emph{Source and licence.} A contact version of Kirby's theorem. Version arXiv:2605.27126v1 (CC BY 4.0), distributed by arXiv under the Creative Commons Attribution 4.0 International licence, http://creativecommons.org/licenses/by/4.0/, as recorded in the arXiv licence metadata for this exact version and shown on the abstract page https://arxiv.org/abs/2605.27126v1. Full licence notice: \texttt{licenses/arxiv-2605.27126-CC-BY-4.0.txt}.\par\smallskip

\noindent\textit{Editorial scope.} Counted unit: the article's lemma lem:Schur\_local (source0.tex lines 2212--2253). The article's complete original proof of it is delivered with it (source0.tex lines 2255--2300, one \texttt{proof} environment of 46 lines). No mathematics is written, repaired or re-derived: the statement and the proof are the article's own lines, in the article's own notation, and no retained line is substituted or re-flowed.  No further source-local context is needed: every cross-reference of every retained line resolves inside the two delivered spans.  Nothing was omitted from any retained span, so the package carries no omission record.  No retained line contains a citation, so the wrapper carries no bibliography and no bibliography entry.  No retained line contains a figure environment, an image-inclusion command or drawing code, so no image is delivered and the package declares no asset dependency.  Editorial formatting only, in addition: an AMS \texttt{amsart} preamble that restates the article's own declarations and macros used by the retained lines, namely the environments \texttt{lemma} on separate counters, together with the standard packages those lines need.  The retained environments print in this document as Lemma 1 in order of appearance, whereas the article prints the same items with the numbering of its own full document.  The complete native source of the article as distributed in the arXiv e-print for this exact version is bundled unchanged as \texttt{sources/201470/source0.tex}.
\begin{lemma}\label{lem:Schur_local}
Let
\[
Q=\begin{pmatrix} A & B^T \\ B & C \end{pmatrix},
\qquad
\boldsymbol r=\binom{\boldsymbol a}{\boldsymbol b},
\qquad
Q'=\begin{pmatrix} A' & (B')^T \\ B' & C \end{pmatrix},
\qquad
\boldsymbol r'=\binom{\boldsymbol a'}{\boldsymbol b},
\]
where \(A\) and \(A'\) are symmetric and invertible\footnote{Here, we use the convention that the empty matrix is invertible and has signature zero. With this convention, Lemma~\ref{lem:Schur_local} also applies when one of the two parts involved in the move is empty.} such that
\[
B A^{-1} B^T = B'(A')^{-1}(B')^T
\qquad\text{and}\qquad
B A^{-1}\boldsymbol a = B'(A')^{-1}\boldsymbol a'.
\]
Then \(Q\boldsymbol x=\boldsymbol r\) admits a rational solution if and only if \(Q'\boldsymbol x'=\boldsymbol r'\) does. In that case, we have
\begin{align*}
    (c')^2-c^2 &= (\boldsymbol a')^T(A')^{-1}\boldsymbol a' - \boldsymbol a^TA^{-1}\boldsymbol a, \textrm{ and }\\
\sigma(Q')-\sigma(Q) &=\sigma(A')-\sigma(A).
\end{align*}
In particular, it follows that if
\[
(\boldsymbol a')^T(A')^{-1}\boldsymbol a'
-
\boldsymbol a^TA^{-1}\boldsymbol a
=
3\bigl(\sigma(A')-\sigma(A)\bigr)
+
2\bigl(\operatorname{dim} A'-\operatorname{dim} A\bigr)
-
4(q'-q),
\]
then the corresponding values of the diagrammatic \(d_3\)-formula agree, i.e.\
\[
d_3^{\mathrm{surg}}(Q',\boldsymbol r',q')
=
d_3^{\mathrm{surg}}(Q,\boldsymbol r,q).
\]
Here \(d_3^{\mathrm{surg}}(Q,\boldsymbol r,q)\) denotes the right-hand side of the surgery formula written in terms of the matrix \(Q\), the rotation vector \(\boldsymbol r\), and the number \(q\) of contact \((+1)\)-surgeries.
\end{lemma}

\begin{proof}
We write \(\boldsymbol x=\binom{\boldsymbol z}{\boldsymbol y}\). Solving the first block equation and substituting into the second shows that \(Q\boldsymbol x=\boldsymbol r\) has a solution $\boldsymbol{x}$ if and only if the equation
\[
(C-BA^{-1}B^T)\boldsymbol y
=
\boldsymbol b-BA^{-1}\boldsymbol a
\]
has a solution $\boldsymbol{y}$; and in that case $\boldsymbol{z}$ is given by $\boldsymbol z=A^{-1}(\boldsymbol a-B^T\boldsymbol y)$.
The same computation shows that \(Q'\boldsymbol x'=\boldsymbol r'\) has a solution if and only if 
\[
(C-B'(A')^{-1}(B')^T)\boldsymbol y
=
\boldsymbol b-B'(A')^{-1}\boldsymbol a',
\]
has a solution. By assumption, these two equations agree, so the solvability of the two systems agrees.
If this is the case, we compute
\[
\boldsymbol r^T\boldsymbol x
=
\boldsymbol a^TA^{-1}\boldsymbol a
+
(\boldsymbol b-BA^{-1}\boldsymbol a)^T\boldsymbol y,
\]
and the same formula holds for \((Q',\boldsymbol r')\). Hence
\[
(\boldsymbol r')^T\boldsymbol x' - \boldsymbol r^T\boldsymbol x
=
(\boldsymbol a')^T(A')^{-1}\boldsymbol a'
-
\boldsymbol a^TA^{-1}\boldsymbol a,
\]
which proves the claimed formula for \((c')^2-c^2\).
Finally, since \(A\) is invertible, \(Q\) is congruent over \(\mathbb R\) to
\[
A\oplus(C-BA^{-1}B^T),
\]
and similarly \(Q'\) is congruent over \(\mathbb R\) to
\[
A'\oplus(C-B'(A')^{-1}(B')^T).
\]
The Schur complements agree by assumption, and therefore
\[
\sigma(Q')-\sigma(Q)=\sigma(A')-\sigma(A).
\]
The last claim follows by substituting these identities into the definition of \(d_3^{\mathrm{surg}}\).
\end{proof}

\end{document}
